\section{Experimental setup}\label{sec:setup}

\paragraph{Hardware and software}
All runs used free Kaggle notebook sessions with two NVIDIA Tesla T4 GPUs (15 GB each), four virtual CPUs and 31.3 GiB of memory. Each process used one GPU; two processes could run in parallel. The code of every session was embedded in its notebook and checked against a fixed version-control tree hash before it ran. Every session recorded its Python version, package versions and model revisions. The readers ran with the same stack in every session (PyTorch 2.13.0 with CUDA 13.0, vLLM 0.30.0, Transformers 5.17.0). Between Phases~1 and~2 the Kaggle image changed from Python 3.12.13 to 3.13.15, and NumPy from 2.0.2 to 2.1.3; Phase~3 used the same image and package versions as Phase~2. The two sessions that encoded the NQ and HotpotQA sub-corpora used older images (PyTorch 2.10.0 and 2.11.0 with CUDA 12.8, Transformers 5.0.0); Table~\ref{tab:env} lists the software of every session as its environment log recorded it.

\paragraph{Determinism across phases}
Some comparisons cross Phases~1 and~2: the Qwen--NQ cell against the new cells in the pooled model, $k=5$ against the new $k$ values, the reranker runs against the Phase~1 runs, and the Phase~3 runs of the Qwen reader against its Phase~1 sub-corpus runs. To check that the software change does not alter the reader, we compared Phases~1 and~2 on inputs that were identical by construction: the gold-only context of Phase~1 and the gold-only condition of the context experiment. Of \ev{CS.n} identical prompts, \ev{CS.same} produced identical outputs; the one difference was a shorter answer to the same question. A rerun of 100 prompts within one session gave \ev{CS.within} identical outputs.

\paragraph{Seeds and repetitions}
The query order (seed 20260926), bootstrap (seed 0), decoding (greedy, seed 0) and the random passages of the context experiment were fixed. Retrieval and greedy decoding are deterministic within a session, so we ran each answer condition once. The cost measurements were repeated five times.

\paragraph{Analysis and verification}
One analysis script reads only the committed raw outputs and writes all statistics, tables and figures. The author reran it on a second computer and obtained identical output. A separate script, written without the project code and using only the Python standard library, re-scored all \ev{IC.n} answers from the raw reader outputs, recomputed every Hit@5 and EM mean, every transfer ratio point estimate, all 84 pairwise EM tests with their Holm adjustment, and Cohen's $\kappa$. It found no disagreement. Bootstrap intervals, F1, answer containment and the pooled model were not re-implemented independently. Phase~3 has its own analysis script, which also reads only committed outputs. It checks that the passages each reader saw equal the top five of the stored ranked lists, and that the sub-corpus intervals it recomputes equal the committed ones. Re-encoding the full corpus gave Hit@5 and nDCG@10 values \ev{V6.e9.repro} to those of the earlier full-corpus check (E9) for all four retrievers. A second script, again written without the project code, re-scored all \ev{V6.n.gen} Phase~3 answers from the raw reader outputs and recomputed the means, ratios with their intervals, McNemar tests, query-level contributions and the answer-string shares of MiniLM-L6 and e5-base-v2; it found no disagreement.

\subsection{Deviations and failed runs}\label{sec:deviations}
We list every departure from the registered design.
\begin{enumerate}
\item \textbf{NQ join gate.} The original design required a join rate of at least 0.70 between BEIR-NQ and NQ-open; the observed rate was \ev{N.nq.joinrate}. Before any outcome was observed, we replaced the gate with a minimum of 1,000 joined queries, because the two resources are independently filtered subsets of Natural Questions. Retrieval results on all \ev{N.nq.judged} judged queries are reported where they matter (Table~\ref{tab:fullcorpus}).
\item \textbf{Software image.} The initial Phase~2 session stopped at the environment check because the Kaggle image had moved to Python 3.13, for which the pinned NumPy had no wheel. We unpinned NumPy and recorded all version differences; no data were produced by the stopped run.
\item \textbf{Failed reader start.} The initial HotpotQA reader session stopped before generating anything because the image's torchaudio build did not match the CUDA version of PyTorch. Later sessions reinstalled the matching build and checked the stack before any other step.
\item \textbf{Reduced scope of E6.} Approximate search was planned for all six encoders. To stay within the GPU quota it ran for MiniLM-L6, bge-small and e5-base-v2, using the sub-corpus rows of the full-corpus embeddings.
\item \textbf{E7 scope.} The repeated cost measurement covered encoding throughput and query-encoding latency; search latency was measured separately in E6 and is not repeated.
\item \textbf{Phase~1 deviations.} We dropped the planned RAGAS cross-check, because its local deployment needed infrastructure that was untested on our hardware. We did not run the planned Matryoshka dimension ablation for Qwen3-Embedding-0.6B, so dimension and model are confounded in our cost comparison. The design planned two annotators for the judge validation; one took part. One gold-passage output reached the token limit, so the judge scored 499 instead of 500 gold items. The initial Phase~1 reading session stopped before generating because a notebook cell ran before the restore step; it was rerun with unchanged code. We estimated the Phase~1 CPU benchmark at one hour; it took eight hours and 24 minutes. This used free quota but changed no data.
\item \textbf{Failure log.} The design promised one machine-readable failure log. Failures were recorded instead in two dated design addenda and in the session logs.
\item \textbf{Phase~3 human audit.} The Phase~3 design adds a blinded audit of 100 seeded e5-base-v2 retrieval misses by the author. Its labels were not complete when this version was written, so it is not reported. The Phase~3 session itself ran as registered, with no failed run.
\end{enumerate}
No run was repeated to obtain a different result, and no result was discarded.
